\documentclass{article}
\usepackage{graphicx} % Required for inserting images
\usepackage{amsmath}
\usepackage{fontspec}
\usepackage{polyglossia}
\setmainlanguage{vietnamese}
\setmainfont{Times New Roman}

\title{University Math Fundamentals in Physics}
\author{Author: Hoang Minh Hieu}
\date{Hanoi, September 2026}

\begin{document}

\section{Vector}
\subsection{Định nghĩa}
\textbf{Vector} có thể hiểu là 1 đoạn thẳng có hướng, đi từ điểm A đến điểm B hoặc từ B đến A, hướng khi đi từ A đến B sẽ ngược lại hướng từ B đến A.\newline
Cho 2 điểm A = ($a_1, a_2,a_3$) và B = ($b_1,b_2,b_3$), \textbf{vector} nối 2 điểm AB sẽ được thể hiện như sau: 
\[\vec{AB}=(b_1-a_1, b_2-a_2,b_3-a_3)\]
Ví dụ: 
\begin{center}
     Cho A = (1, 2, 3), B = (4, 5, 6), $\vec{AB}$ = (4-1, 5-2, 6-3) = (3, 3, 3)
\end{center}
\underline{\textbf{Vector} đơn vị} là các \textbf{vector} có độ dài là 1, đi theo các hướng Ox, Oy, Oz. Với $\hat{i}$ = (1,0,0); $\hat{j}$ = (0,1,0); $\hat{k}$ = (0,0,1).

\subsection{Nhân 1 số với \textbf{vector}}
Khi ta nhân 1 số thực với 1 \textbf{vector}, ta đang kéo giản \textbf{vector} đó với số lần số thực đó. Có thể hiểu là giá trị của số thực đó là bội số của độ dài \textbf{vector} đó. \newline
Ví dụ: \newline
2$\cdot\vec{a}$: Kéo dài \textbf{vector} a ra 2 lần\newline
1/2$\cdot\vec{b}$: Kéo dài \textbf{vecotr} b còn 1/2 lần (hay còn gọi là nén)\newline
-1$\cdot\vec{c}$: Dấu trừ sẽ đảo chiều \textbf{vector}, nếu không có thêm bội số thì sẽ không làm thay đổi chiều dài \textbf{vector}

\subsection{Cộng, trừ 2 vector}
Khi cộng 2 \textbf{vector} ta cộng lại các toạ độ x, y, z của 2 \textbf{vector} mà ta tính.\newline
Ví dụ: 
\[\vec{a}=(1,5,-3),\vec{b}=(-2,-4,5),\vec{a}+\vec{b}=(1-2,5-4,-3+5)=(-1,1,2)\]
Khi trừ 2 \textbf{vector} ta làm tương tự với phép cộng nhưng thay bằng dấu trừ.\newline
Ví dụ: 
\[\vec{c}=(2,3,5),\vec{d}=(5,-1,-3, \vec{c}-\vec{d}=(2-5,3-(-1),5-(-3))=(-3,4,8)\]

\subsection{Tích vô hướng, có hướng}
Đối với tích vô hướng giữa 2 \textbf{vector}, ta nhân từng phần tử x, y, z tương ứng của chúng lại với nhau và cộng tổng tất cả lại, kết quả sẽ là 1 vô hướng (1 số thực).\newline
Ví dụ: 
\[\vec{a}=(0,3,7), \vec{b}=(1,3, 1/2),\vec{a}\cdot\vec{b}=0\cdot1+3\cdot3+7\cdot1/2=25/2\]
Khi tính tích có hướng, ta tính định thức của 1 ma trận 3$\times$3 với 2 \textbf{vector} được đặt theo hàng ngang và ở hàng thứ nhất sẽ là các \textbf{vector} đơn vị $\hat{i}$, $\hat{j}$, $\hat{k}$.
Cho 2 \textbf{vector} $\vec{a}=(a_1, b_1, c_1)$ và $\vec{b}=(a_2, b_2, c_2)$: 
\[\vec{a}\times\vec{b}= \begin{vmatrix}
    \hat{i}&\hat{j}&\hat{k}\\
    a_1&b_1&c_1\\
    a_2&b_2&c_2
\end{vmatrix}=\hat{i}\cdot\begin{vmatrix}
    b_1&c_1\\
    b_2&c_2
\end{vmatrix}-\hat{j}\cdot\begin{vmatrix}
    a_1&c_1\\
    a_2&c_2
\end{vmatrix}+\hat{k}\cdot\begin{vmatrix}
    a_1&b_1\\
    a_2&b_2
\end{vmatrix}\]

\section{Giải tích vector}
\subsection{Hàm vector}
1 hàm \textbf{vector} được định nghĩa là: 
\[\vec{F} = \langle P, Q, R \rangle=\hat{i}\cdot P+\hat{j}\cdot Q+\hat{k}\cdot R\]
Với P, Q, R là các hàm số riêng biệt.

\subsection{Giới hạn của hàm vector}
Giới hạn của hàm \textbf{vector} là giới hạn của từng hàm bên trong hàm \textbf{vector} đó.
\[\lim_{x \to a}\vec{F}= \langle \lim_{x \to a}P, \lim_{x \to a}Q,\lim_{x \to a}R \rangle\]
Ví dụ: 
Tính giới hạn của $\vec{F} = \langle 2x^2, ln(3x), \frac{1}{2} \rangle$ khi x tiến tới $\infty$.
\[lim_{x \to \infty}\vec{F}=\langle lim_{x \to \infty}2x^2, lim_{x \to \infty}ln(3x), lim_{x \to \infty}1/2 \rangle\]
\[lim_{x \to \infty}\vec{F} = \langle \infty, \infty, \frac{1}{2} \rangle\]

\subsection{Đạo hàm của hàm vector}
Đạo hàm của hàm \textbf{vector} là đạo hàm của từng hàm bên trong hàm \textbf{vector} đó.
\[\frac{d \vec{F}}{dx}=\langle \frac{dP}{dx}, \frac{dQ}{dx}, \frac{dR}{dx} \rangle\]
Ví dụ: Tính đạo hàm của $\vec{F} = \langle 2x^2, ln(3x), \frac{1}{2} \rangle$.
\[\frac{d \vec{F}}{dx}=\langle \frac{d}{dx}2x^2, \frac{d}{dx}ln(3x), \frac{d}{dx}\frac{1}{2} \rangle\]
\[\frac{d \vec{F}}{dx}=\langle 4x, \frac{1}{x}, 0 \rangle\]

\subsection{Nguyên hàm của hàm vector}
Nguyên hàm của hàm \textbf{vector} là nguyên hàm của từng hàm bên trong hàm \textbf{vector} đó.
\[\int \vec{F}dx = \langle \int Pdx, \int Qdx,\int Rdx\rangle\]
Ví dụ: Tính nguyên hàm của $\vec{F} = \langle 4x, \frac{1}{x}, -sinx \rangle$.
\[\int \vec{F}dx = \langle \int 4xdx, \int \frac{1}{x}dx,\int -sinxdx\rangle\]
\[\int \vec{F}dx = \langle 2x^2, ln(x), cosx \rangle\]

\subsection{Đạo hàm riêng}
Đạo hàm riêng là đạo hàm theo 1 biến của hàm số nhiều biến. Ký hiệu là $\frac{\partial f}{\partial x}, \frac{\partial f}{\partial y}$. Đạo hàm theo x hoặc y thì sẽ coi toàn bộ tất cả biến còn lại như là 1 hằng số.
\begin{center}
    Cho f(x) = 2xy + 2x$y^2$ -3, $\frac{\partial f}{\partial x} = 2y + 2y^2, \frac{\partial f}{\partial y}=2x+4xy$.
\end{center}
Đạo hàm riêng là 1 mắt chốt trong Divergence và Curl(hay còn gọi là Rot).

\subsection{Gradient ($\nabla$)}
Gradient (hoặc Nabla) được định nghĩa như sau: 
\[\nabla = \hat{i}\cdot\frac{\partial}{\partial x}+\hat{j}\cdot\frac{\partial}{\partial y}+\hat{k}\cdot\frac{\partial}{\partial z}\]
Tích vô hướng giữa Gradient và hàm \textbf{vector}, được gọi là \textbf{Divergence}: 
\[\nabla\cdot \vec{F} = (\frac{\partial}{\partial x} + \frac{\partial}{\partial y} + \frac{\partial}{\partial z})\cdot\begin{pmatrix}
    P\\Q\\R
\end{pmatrix}=\frac{\partial P}{\partial x}+\frac{\partial Q}{\partial y}+\frac{\partial R}{\partial z}\]
Kết quả của phép tích vô hướng giữa Gradient và hàm \textbf{vector} sẽ là 1 số thực.\newline
Tích có hướng giữa Gradient và hàm \textbf{vector}, được gọi là \textbf{Curl} hoặc \textbf{Rot}: 
\[\nabla \times \vec{F}=\begin{vmatrix}
    \hat{i}&\hat{j}&\hat{k}\\
    \frac{\partial}{\partial x}&\frac{\partial}{\partial y}&\frac{\partial}{\partial z}\\
    P&Q&R
\end{vmatrix}=\hat{i}\cdot\begin{vmatrix}
    \frac{\partial}{\partial y}&\frac{\partial}{\partial z}\\
    Q&R
\end{vmatrix}-\hat{j}\cdot\begin{vmatrix}
    \frac{\partial}{\partial x}& \frac{\partial}{\partial z}\\
    P&R
\end{vmatrix}+\hat{k}\cdot\begin{vmatrix}
    \frac{\partial}{\partial x}&\frac{\partial}{\partial y}\\
    P&Q
\end{vmatrix}\]
Kết quả của phép tích có hướng giữa Gradient và hàm \textbf{vector} sẽ là 1 \textbf{}{vector}.

\end{document}
